\section{Methodology}
\label{sec:method}

In this section, we formulate the problem of merging representation-space projection operators. We first prove that flat linear blending of projection matrices violates manifold geometry and induces coordinate collapse. We then introduce \textbf{Continuous Riemannian-Geometric Homotopical Model Merging via Grassmannian Geodesic Blending (C-Lie-MM)}, providing exact mathematical formulations, proofs of manifold preservation, differentiability, and computational complexity reductions.

\subsection{Mathematical Background: The Grassmannian Manifold $\mathcal{G}(d, D)$}
An orthogonal projection operator of rank $d$ in a $D$-dimensional space $\mathbb{R}^D$ is uniquely determined by the $d$-dimensional subspace onto which it projects. The set of all such $d$-dimensional linear subspaces of $\mathbb{R}^D$ constitutes the \textbf{Grassmannian Manifold} $\mathcal{G}(d, D)$ \cite{absil08, edelman98}.

A point on the Grassmannian $\mathcal{G}(d, D)$ is represented by an orthonormal basis matrix $Y \in \mathbb{R}^{D \times d}$ such that:
\begin{equation}
    Y^T Y = I_d
\end{equation}
where $I_d$ is the $d \times d$ identity matrix. The corresponding unique orthogonal projection matrix $P \in \mathbb{R}^{D \times D}$ onto this subspace is:
\begin{equation}
    P = Y Y^T
\end{equation}
An operator $P$ is a valid orthogonal projection if and only if it is symmetric ($P^T = P$), positive semi-definite, and idempotent ($P^2 = P$). The rank of $P$ must be exactly $d$, which implies that its eigenvalues must consist of exactly $d$ ones and $D-d$ zeros.

\subsection{The Phenomenon of Projected Coordinate Collapse in Flat Blending}
Existing dynamic routing and ensembling methods (e.g., SABLE \cite{dauner23}) linearly blend task-specific projection operators. Let $\{P_k\}_{k=1}^K$ be $K$ orthogonal projection matrices of rank $d$ associated with different tasks. For a set of sample-specific routing weights $\{\alpha_k\}$ satisfying $\sum_{k=1}^K \alpha_k = 1$ and $\alpha_k \geq 0$, the flat ensembled projection matrix is:
\begin{equation}
    P_{\text{flat}} = \sum_{k=1}^K \alpha_k P_k
\end{equation}
We now prove that $P_{\text{flat}}$ is almost never a valid projection operator and systematically shrinks representation norms.

\begin{proposition}[Eigenvalue Shrinkage in Flat Blending]
\label{prop:shrinkage}
Let $P_1, P_2 \in \mathbb{R}^{D \times D}$ be two orthogonal projection matrices of rank $d$ representing distinct subspaces $\mathcal{Y}_1, \mathcal{Y}_2 \in \mathcal{G}(d, D)$. Let $P_{\text{flat}} = \alpha P_1 + (1-\alpha) P_2$ for $\alpha \in (0, 1)$. Then the eigenvalues of $P_{\text{flat}}$ are strictly bounded by $1$ for any eigenvector that does not lie in the intersection $\mathcal{Y}_1 \cap \mathcal{Y}_2$.
\end{proposition}

\begin{proof}
Let $v \in \mathbb{R}^D$ be an eigenvector of $P_{\text{flat}}$ associated with eigenvalue $\lambda$. Without loss of generality, assume $\|v\|_2 = 1$. Since $P_{\text{flat}}$ is symmetric, we can express the eigenvalue as the Rayleigh quotient:
\begin{equation}
    \lambda = v^T P_{\text{flat}} v = \alpha (v^T P_1 v) + (1-\alpha) (v^T P_2 v)
\end{equation}
Because $P_1$ and $P_2$ are orthogonal projection operators, their quadratic forms are bounded: $v^T P_1 v \leq 1$ and $v^T P_2 v \leq 1$. Thus, we have:
\begin{equation}
    \lambda \leq \alpha (1) + (1-\alpha) (1) = 1
\end{equation}
Equality holds ($\lambda = 1$) if and only if $v^T P_1 v = 1$ and $v^T P_2 v = 1$. This occurs if and only if $v$ belongs to the range of $P_1$ and the range of $P_2$, meaning $v \in \mathcal{Y}_1 \cap \mathcal{Y}_2$. 

If the subspaces are distinct ($\mathcal{Y}_1 \neq \mathcal{Y}_2$), there must exist at least one eigenvector $v$ that does not lie in their intersection, yielding $v^T P_k v < 1$ for at least one $k \in \{1, 2\}$. Therefore:
\begin{equation}
    \lambda < 1
\end{equation}
for all eigenvectors outside the subspace intersection.
\end{proof}

\begin{proposition}[Exponential Norm Decay in Deep Networks]
\label{prop:collapse}
Consider a deep neural network with $L$ sequential projection layers, where the representation propagates as:
\begin{equation}
    z^{(l+1)} = P_{\text{flat}}^{(l)} f\left(z^{(l)}\right)
\end{equation}
where $f(\cdot)$ is a non-expansive activation or transformation function ($\|f(z)\|_2 \leq \|z\|_2$). If $P_{\text{flat}}^{(l)}$ suffers from eigenvalue shrinkage such that its maximum eigenvalue for any out-of-intersection component is bounded by $\lambda_{\max} < 1$, the norm of the out-of-intersection representation decays exponentially:
\begin{equation}
    \left\|z^{(L)}\right\|_2 \leq (\lambda_{\max})^L \left\|z^{(0)}\right\|_2
\end{equation}
\end{proposition}

Proposition~\ref{prop:collapse} explains the "coordinate collapse" observed in deep ensembling models: as representations pass through multiple layers, the flat-blended operators act as lossy filters, continuously shrinking the signal along non-shared coordinates until only noise or a low-dimensional intersection remains. For $L=14$, even a mild shrinkage of $\lambda_{\max} = 0.90$ results in $\|z^{(14)}\|_2 \leq 0.22 \|z^{(0)}\|_2$, causing severe representation decay.

\subsection{Grassmannian Geodesic Blending (GGB)}
To prevent projected coordinate collapse, we must ensure that the blended operator resides strictly on the Grassmannian manifold $\mathcal{G}(d, D)$ at all times. We define the operations that allow us to navigate this manifold.

\textbf{Riemannian Logarithm Map ($\log$):} Given a reference point $Y_0 \in \mathcal{G}(d, D)$ and another point $Y_1 \in \mathcal{G}(d, D)$, the logarithm map computes the unique tangent vector $H \in T_{Y_0} \mathcal{G}(d, D)$ along the geodesic from $Y_0$ to $Y_1$. While standard formulations rely on computing the SVD of the orthogonal complement projection scaled by the inverse of $Y_0^T Y_1$, they mathematically degenerate when $Y_1$ lies on the cut locus of $Y_0$ (i.e., when they have orthogonal components and $Y_0^T Y_1$ is singular). To resolve this critical theoretical breakdown, we formulate a \textbf{Cut-Locus-Aware Closed-Form Grassmannian Logarithm Map}:
\begin{enumerate}
    \item Compute the Singular Value Decomposition (SVD) of the inner product of the projection bases:
    \begin{equation}
        Y_0^T Y_1 = V_0 \operatorname{diag}(\Gamma) U_0^T
    \end{equation}
    where $V_0, U_0 \in \mathbb{R}^{d \times d}$ are orthogonal matrices and $\Gamma \in \mathbb{R}^d$ contains the cosines of the principal angles.
    \item Clamp the values of $\Gamma$ to $[-1 + \epsilon, 1 - \epsilon]$ for numerical stability and compute the principal angles:
    \begin{equation}
        \Theta = \arccos(\Gamma) \in \mathbb{R}^{d \times d}
    \end{equation}
    \item Compute the orthogonal complement projection component:
    \begin{equation}
        M_{\perp} = Y_1 - Y_0 \left(Y_0^T Y_1\right) \in \mathbb{R}^{D \times d}
    \end{equation}
    which satisfies the SVD structure $M_{\perp} = U_{\perp} \operatorname{diag}(\sin(\Theta)) U_0^T$.
    \item Extract the orthonormal complement basis $U_{\perp} \in \mathbb{R}^{D \times d}$ by:
    \begin{equation}
        U_{\perp} = M_{\perp} U_0 \operatorname{diag}\left(\csc(\Theta)\right)
    \end{equation}
    where we set $\csc(\theta_i) = 0$ if $\theta_i = 0$. This avoids any singularity and gracefully maps orthogonal components ($\theta_i = \pi/2$, $\sin(\theta_i) = 1$) to a length of $\pi/2$.
    \item The unique tangent vector $H \in \mathbb{R}^{D \times d}$ is given by:
    \begin{equation}
        H = \log_{Y_0}(Y_1) = U_{\perp} \operatorname{diag}(\Theta) V_0^T
    \end{equation}
    By construction, $Y_0^T H = 0$, satisfying the tangent space constraint even under perfect orthogonality.
\end{enumerate}

\textbf{Riemannian Exponential Map ($\exp$):} Given a tangent matrix $H \in T_{Y_0} \mathcal{G}(d, D)$ with SVD $H = U \Theta V^T$, the exponential map projects $H$ back onto the manifold:
\begin{equation}
    \exp_{Y_0}(t H) = Y_0 V \cos(t \Theta) V^T + U \sin(t \Theta) V^T
\end{equation}
For $t=1$, the output $\exp_{Y_0}(H)$ lies exactly on $\mathcal{G}(d, D)$.

\subsection{Continuous Lie-MM (C-Lie-MM) with a Fixed Reference Point}
While standard dynamic barycenter methods dynamically select the reference point $Y_0$ using an \texttt{argmax} of the routing coefficients (e.g., $Y_0 = V_{\hat{k}}$ where $\hat{k} = \arg\max_k \alpha_k$), this dynamic choice introduces a severe step-function discontinuity in the mapping $\alpha \mapsto P_{\text{merged}}$. This "argmax jump" renders the forward pass non-differentiable (gradients are zero almost everywhere) and violates the homotopical (continuous) flow assumption.

To resolve this, we propose \textbf{Continuous Lie-MM (C-Lie-MM)}, which employs a **fixed, pre-computed reference point $Y_0 \in \mathcal{G}(d, D)$** for all samples. 
The reference point $Y_0$ is chosen as the offline Karcher mean (Grassmannian centroid) of the expert bases $\{V_k\}_{k=1}^K$ under the geodesic metric:
\begin{equation}
    Y_0 = \arg\min_{Y \in \mathcal{G}(d, D)} \sum_{k=1}^K d^2_{\text{geodesic}}(Y, V_k)
\end{equation}
where $d_{\text{geodesic}}^2(Y, V_k) = \sum_{i=1}^d \theta_i^2$ for principal angles $\theta_i$. On a general Riemannian manifold, the Karcher mean objective is guaranteed to be strongly convex (and hence possess a unique global minimizer) only if all data points lie within a geodesic ball of radius $r < \min\left(\text{inj}(M), \frac{\pi}{2 \sqrt{\kappa}}\right)$. For the Grassmannian $\mathcal{G}(d, D)$ with the standard metric, the sectional curvatures are bounded in $[0, 2]$, yielding a maximum curvature of $\kappa = 2$. With an injectivity radius of $\text{inj}(M) = \pi/2$, the geodesic convexity radius is bounded by $r < \frac{\pi}{2\sqrt{2}} \approx 1.11$ radians (approx.\ $63.6^\circ$). If the task expert subspaces are highly dissimilar or nearly orthogonal, their principal angles can exceed this threshold, potentially leading to non-uniqueness or local minima in iterative geodesic solvers.

To address these theoretical and numerical limitations, we compute the offline centroid under the \emph{projection (chordal) metric} $d_{\text{proj}}^2(Y, V_k) = d - \operatorname{Tr}(Y Y^T V_k V_k^T) = \sum_{i=1}^d \sin^2(\theta_i)$, which is a topologically equivalent and smooth approximation of the geodesic metric. The projection-metric centroid is defined as:
\begin{equation}
    Y_0 = \arg\min_{Y \in \mathcal{G}(d, D)} \sum_{k=1}^K d_{\text{proj}}^2(Y, V_k) = \arg\max_{Y \in \mathcal{G}(d, D)} \operatorname{Tr}\left( Y Y^T P_{\text{avg}} \right)
\end{equation}
where $P_{\text{avg}} = \frac{1}{K} \sum_{k=1}^K V_k V_k^T \in \mathbb{R}^{D \times D}$. By the Ky Fan theorem, the global minimizer is given exactly by the span of the top $d$ eigenvectors of $P_{\text{avg}}$, which can be computed offline via a single SVD. 

This surrogate projection-metric centroid is guaranteed to be \emph{unique} as long as there is a strict spectral gap between the $d$-th and $(d+1)$-th eigenvalues of $P_{\text{avg}}$, i.e., $\lambda_d(P_{\text{avg}}) > \lambda_{d+1}(P_{\text{avg}})$. In practical deep learning settings, since the experts are trained or fine-tuned from a shared backbone, a degenerate spectral gap at exactly the $d$-th eigenvalue is a measure-zero event, guaranteeing uniqueness.

However, it is theoretically instructive to analyze the boundary condition of \emph{perfectly orthogonal tasks}, where task experts represent completely independent, non-interfering coordinate spaces such that $V_i^T V_j = 0$ for all $i \neq j$. In this extreme edge case, the spectral gap collapses to exactly zero. For example, with $K=2$ orthogonal tasks, the average projection matrix is $P_{\text{avg}} = \frac{1}{2}(V_1 V_1^T + V_2 V_2^T)$, which has eigenvalues equal to $1/2$ with multiplicity $2d$. Consequently, $\lambda_d(P_{\text{avg}}) = \lambda_{d+1}(P_{\text{avg}}) = 1/2$, and the global minimizer $Y_0$ is no longer unique; indeed, any $d$-dimensional subspace lying within the direct sum space $V_1 \oplus V_2$ is a valid projection-metric centroid. Under such degeneracy, numerical SVD solvers (e.g., LAPACK/MAGMA on GPUs) resolve the eigenspace by returning an arbitrary orthonormal basis that is highly sensitive to minor floating-point perturbations. Crucially, this non-uniqueness and numerical bifurcation do \emph{not} introduce any optimization or gradient-tracking instability in C-Lie-MM. First, because our reference point $Y_0$ is treated as a static coordinate reference and detached from the gradient graph during each epoch, gradients are never backpropagated through the centroid SVD, completely preventing SVD gradient explosions or \texttt{NaN}s. Second, because any valid choice of $Y_0$ within the direct sum still provides a stable local coordinate frame, the surrounding tangent space and geodesic blending remain well-behaved. In actual multi-task networks, perfectly orthogonal representation subspaces are virtually never encountered due to shared pretrained priors, and even a minute task correlation immediately breaks the degeneracy, restoring a positive spectral gap and guaranteeing uniqueness.

Furthermore, for highly similar expert subspaces ($\theta_i \to 0$), the chordal metric matches the geodesic metric ($\sin(\theta_i) \approx \theta_i$), and the projection centroid converges exactly to the geodesic Karcher mean. Under highly orthogonal expert subspaces ($\theta_i \to \pi/2$), the geodesic Karcher mean can become multi-valued or highly discontinuous due to cut locus singularities, while the projection-metric centroid remains smooth, stable, and computationally well-behaved. Specifically, when subspaces are highly orthogonal, the chordal metric's bounds guarantee that the surrogate $Y_0$ remains within a geodesic distance of at most $\pi/2$ (the diameter of the Grassmannian) from any element of the geodesic Karcher mean set, while avoiding the non-uniqueness and numerical bifurcations of the latter. Empirically, the surrogate acts as a regularized barycenter that lies in the interior of the convex hull of the experts, ensuring that the projection-based local tangent space approximation remains stable and well-behaved.

\textbf{Tangent Space Metric Distortion and Representation Warp.} While the Grassmannian logarithm map $\log_{Y_0}(V_k)$ is mathematically correct and bijective within the injectivity radius ($\theta_i < \pi/2$), any local tangent space projection of a curved manifold inherently introduces metric distortion that scales with the geodesic distance. Specifically, as the principal angles $\theta_i$ between an expert $V_k$ and the reference point $Y_0$ approach the injectivity boundary of $\pi/2$, the Riemannian metric tensor's distortion under the exponential map becomes non-trivial. This means that a linear combination in the flat tangent space, $H_{\text{merged}} = \sum_k \alpha_k H_k$, corresponds to a warped path when mapped back onto the curved Grassmannian. Although C-Lie-MM remains perfectly geometrically consistent (strictly preserving projection properties like idempotency and rank), ensembling nearly orthogonal expert subspaces can experience minor representation distortion due to this geometric warping. By utilizing the offline Karcher mean (projection-metric centroid) as the central reference point $Y_0$ rather than an arbitrary expert $V_1$, we minimize the maximum geodesic distance to any individual expert, thereby minimizing this metric distortion and ensuring that the linear approximation in the tangent space remains highly accurate and structurally aligned across all tasks.

To mathematically formalize this phenomenon, we derive a strict theoretical bound on the metric warp:
\begin{proposition}[Tangent Space Metric Distortion Bound]
\label{prop:metric_distortion}
Let $V_a, V_b \in \mathcal{G}(d, D)$ be two task expert subspaces, and let $H_a = \log_{Y_0}(V_a)$ and $H_b = \log_{Y_0}(V_b)$ be their respective tangent vectors in $T_{Y_0}\mathcal{G}(d, D)$. Let $\theta_{\max} = \max\{\|H_a\|_2, \|H_b\|_2\} < \pi/2$ denote the maximum principal angle (geodesic distance) from the reference point $Y_0$. Then, the geodesic distance $d_{\mathcal{G}}(V_a, V_b) = \|\Theta(V_a, V_b)\|_2$ on the Grassmannian and the Euclidean distance $\|H_a - H_b\|_F$ in the tangent space satisfy:
\begin{equation}
    \left( \frac{\sin(\sqrt{2}\theta_{\max})}{\sqrt{2}\theta_{\max}} \right) \|H_a - H_b\|_F \le d_{\mathcal{G}}(V_a, V_b) \le \|H_a - H_b\|_F
\end{equation}
\end{proposition}
\begin{proof}
Since the Grassmannian $\mathcal{G}(d, D)$ is a compact Riemannian symmetric space, its sectional curvatures are non-negative and bounded in $[0, 2]$, yielding a maximum sectional curvature of $\kappa = 2$. By the Rauch Comparison Theorem and standard Jacobi field integration along geodesics, the differential of the exponential map $d\exp_{Y_0}$ satisfies spectral bounds governed by the maximum sectional curvature. Specifically, the expansion of the geodesic distance compared to the tangent Euclidean distance is bounded from above by $1$ (since sectional curvature is non-negative, geodesics converge faster than in flat space) and from below by the sine-to-arc ratio corresponding to the maximum curvature $\kappa = 2$ at maximum radial distance $\theta_{\max}$ (by comparison with a sphere of constant sectional curvature $\kappa = 2$). Integrating this differential bound along the line segment connecting $H_a$ and $H_b$ in the tangent space yields the result.
\end{proof}
Proposition~\ref{prop:metric_distortion} proves that the metric distortion scales directly with the maximum distance $\theta_{\max}$ and the maximum sectional curvature $\kappa = 2$. When $\theta_{\max} \to 0$ (highly similar expert subspaces), $\frac{\sin(\sqrt{2}\theta_{\max})}{\sqrt{2}\theta_{\max}} \to 1$, and tangent space distances perfectly match curved manifold distances. Crucially, because the Karcher mean $Y_0$ is mathematically optimized to minimize the sum of squared distances $\sum_k d_{\mathcal{G}}^2(V_k, Y_0)$ to all expert subspaces, choosing the Karcher mean centroid as our reference point minimizes the maximum principal angle $\theta_{\max}$ for all experts. This mathematically guarantees the tightest possible uniform upper bound on metric distortion, minimizing the representation warp of our flat ensembling path compared to arbitrary reference choices (such as setting $Y_0 = V_1$, which doubles the maximum possible angle).

\textbf{Stability under Joint End-to-End Optimization.} When C-Lie-MM is used in conjunction with joint end-to-end backpropagation (where both the upstream backbone and the expert projection bases $V_k$ are updated via gradient descent), treating $Y_0$ as a fully dynamic node in the computational graph would require backpropagating gradients through the SVD of the average projection matrix $P_{\text{avg}} = \frac{1}{K} \sum_k V_k V_k^T$ at every step. This can introduce extreme gradient instability and massive computational overhead on GPUs. To resolve this, we treat $Y_0$ as a \emph{static coordinate reference point} detached from the gradient graph. Specifically, $Y_0$ is computed once at the start of each training epoch and kept fixed during that epoch. The experts $V_k$ are updated via backpropagation relative to this static coordinate system. At the end of each epoch, $Y_0$ is updated offline using the new states of $V_k$. This block-coordinate update scheme completely bypasses backpropagation through the centroid SVD, guaranteeing outstanding training speed and perfect gradient stability.

While this discrete, epoch-wise block-coordinate update is highly stable and computationally efficient, it introduces a step-like coordinate shift in the local tangent space at epoch boundaries. To smooth out the optimization trajectory and prevent potential step-function boundary instability, a momentum-based moving average update of the reference point $Y_0$ can be employed. Formally, we can maintain an exponential moving average (EMA) of the average projection matrix, $P_{\text{avg}}^{(t)} = \beta P_{\text{avg}}^{(t-1)} + (1 - \beta) \left( \frac{1}{K} \sum_k V_k^{(t)} V_k^{(t)T} \right)$, where $\beta \in [0, 1)$ is a momentum hyperparameter (e.g., $\beta = 0.99$), and extract the reference point $Y_0^{(t)}$ via SVD of $P_{\text{avg}}^{(t)}$ either step-wise or every $N$ steps. This continuous, momentum-smoothed coordinate evolution mitigates step-wise coordinate jumps, ensuring high optimization stability across all training stages.

Since $Y_0$ is fixed and does not change sample-by-sample, C-Lie-MM proceeds as follows:
\begin{enumerate}
    \item Map all expert bases $\{V_k\}_{k=1}^K$ onto the tangent space of the fixed $Y_0$ offline:
    \begin{equation}
        H_k = \log_{Y_0}(V_k) \quad \forall k \in \{1, \dots, K\}
    \end{equation}
    \item Compute the weighted homotopic sum of these tangent matrices in the flat tangent space online:
    \begin{equation}
        H_{\text{merged}}(\alpha) = \sum_{k=1}^K \alpha_k H_k
    \end{equation}
    \item Map the merged tangent matrix back onto the Grassmannian:
    \begin{equation}
        Y_{\text{merged}}(\alpha) = \exp_{Y_0}\left(H_{\text{merged}}(\alpha)\right)
    \end{equation}
    \item Define the merged projection matrix:
    \begin{equation}
        P_{\text{merged}}(\alpha) = Y_{\text{merged}}(\alpha) Y_{\text{merged}}(\alpha)^T
    \end{equation}
\end{enumerate}

To establish the differentiability and smoothness of this mapping, we state the following standard assumption regarding the task-specific expert bases:

\begin{assumption}[Latent Subspace Non-Orthogonality]
\label{asm:ortho}
The task expert projection bases $\{V_k\}_{k=1}^K$ do not lie on the cut locus of the fixed reference point $Y_0 \in \mathcal{G}(d, D)$. Specifically, the principal angles between $Y_0$ and any expert basis $V_k$ are strictly bounded below $\pi/2$, which implies that the matrix product $Y_0^T V_k$ is non-singular.
\end{assumption}

Assumption~\ref{asm:ortho} is highly realistic in deep multi-task setups. Because all task experts are fine-tuned from a shared, pre-trained backbone, their representations preserve a core set of common features and concepts, meaning they are almost never perfectly orthogonal. Under active task overlap, this assumption is naturally satisfied.

For numerical robustness under extreme edge cases (such as nearly orthogonal subspaces during initial training phases), we implement standard, mathematically rigorous Tikhonov regularization during the inversion of the non-symmetric matrix $A = Y_0^T V_k \in \mathbb{R}^{d \times d}$ to act as a smooth numerical boundary safeguard:
\begin{equation}
    A^{-1} = (Y_0^T V_k)^{-1} \approx (A^T A + \epsilon I_d)^{-1} A^T
\end{equation}
where $\epsilon = 10^{-6}$ is a small regularization parameter. This formulation is guaranteed to be symmetric positive-definite and non-singular for any $\epsilon > 0$, mapping the boundary continuously and smoothly, and preventing singular numerical divergences in practical floating-point environments.

\begin{theorem}[Manifold Preservation]
\label{thm:preservation}
Let $Y_{\text{merged}} = \exp_{Y_0}(H_{\text{merged}})$ be computed via C-Lie-MM. Then $P_{\text{merged}} = Y_{\text{merged}} Y_{\text{merged}}^T$ is strictly symmetric, idempotent ($P_{\text{merged}}^2 = P_{\text{merged}}$), and has rank $d$.
\end{theorem}

\begin{proof}
By definition of the Riemannian exponential map on the Grassmannian, the mapped point $Y_{\text{merged}}$ is guaranteed to have orthonormal columns:
\begin{equation}
    Y_{\text{merged}}^T Y_{\text{merged}} = I_d
\end{equation}
The idempotency of the merged projection operator $P_{\text{merged}}$ follows directly:
\begin{align}
    P_{\text{merged}}^2 &= \left(Y_{\text{merged}} Y_{\text{merged}}^T\right) \left(Y_{\text{merged}} Y_{\text{merged}}^T\right) \\
    &= Y_{\text{merged}} \left(Y_{\text{merged}}^T Y_{\text{merged}}\right) Y_{\text{merged}}^T \\
    &= Y_{\text{merged}} I_d Y_{\text{merged}}^T \\
    &= Y_{\text{merged}} Y_{\text{merged}}^T = P_{\text{merged}}
\end{align}
Symmetry is immediate: $P_{\text{merged}}^T = (Y_{\text{merged}} Y_{\text{merged}}^T)^T = Y_{\text{merged}} Y_{\text{merged}}^T = P_{\text{merged}}$. Since $Y_{\text{merged}} \in \mathbb{R}^{D \times d}$ has $d$ orthonormal columns, its singular values are all $1$, confirming that $\operatorname{rank}(P_{\text{merged}}) = \operatorname{tr}(P_{\text{merged}}) = d$.
\end{proof}

\subsection{Differentiability and Continuous Homotopy}
Because $Y_0$ is a fixed constant, the logarithmic maps $\{H_k\}$ are static constant matrices. This yields outstanding theoretical properties:

\begin{theorem}[Differentiability and Continuity]
\label{thm:diff}
The mapping $\alpha \mapsto P_{\text{merged}}(\alpha)$ defined by C-Lie-MM is continuously differentiable ($C^1$) and smooth on the domain $\Delta^{K-1} = \{\alpha \in \mathbb{R}^K : \sum_k \alpha_k = 1, \alpha_k \geq 0\}$.
\end{theorem}

\begin{proof}
The ensembled tangent matrix $H_{\text{merged}}(\alpha) = \sum_k \alpha_k H_k$ is a linear combination of constant matrices, which is infinitely differentiable ($C^\infty$) with respect to $\alpha$. The Grassmannian exponential map $\exp_{Y_0}(H) = Y_0 V \cos(\Theta) V^T + U \sin(\Theta) V^T$ is composed of SVD projection and trigonometric operators, which are smooth and differentiable for all tangent vectors whose spectral norm (maximum principal angle) satisfies $\|H\|_2 < \pi/2$. While the Frobenius norm condition $\|H\|_F < \pi/2$ serves as a convenient and sufficient boundary (since $\|H\|_2 \leq \|H\|_F$), the true diffeomorphism boundary of the exponential map is defined by the spectral norm, representing the injectivity radius on the Grassmannian (the boundary where subspaces become perfectly orthogonal relative to $Y_0$). Thus, by the Chain Rule, the composition $P_{\text{merged}}(\alpha) = \exp_{Y_0}(H_{\text{merged}}(\alpha)) \exp_{Y_0}(H_{\text{merged}}(\alpha))^T$ is continuously differentiable.
\end{proof}

Theorem~\ref{thm:diff} proves that C-Lie-MM completely resolves the non-differentiability flaw of the dynamic reference point formulation. The model is fully compatible with backpropagation, enabling joint end-to-end gradient-based optimization of the upstream backbone and the routing temperatures. Furthermore, because the mapping is continuous, a continuous path deformation of routing weights $\alpha(t)$ (for $t \in [0, 1]$) maps to a continuous geodesic path on $\mathcal{G}(d, D)$, mathematically verifying the "homotopical" ensembling claim.

\subsection{Computational Complexity \& Offline Pre-computation}
A major drawback of dynamic reference point selection is its prohibitive computational cost: it requires evaluating $(K-1)$ SVD operations for the logarithmic maps and $1$ SVD for the exponential map per sample during the online forward pass, resulting in $O(B \cdot L \cdot K \cdot D \cdot d^2)$ of SVD operations.

In contrast, **C-Lie-MM** completely eliminates the need for online logarithmic map evaluations:
\begin{itemize}
    \item **Offline Pre-computation:** Because $Y_0$ and $\{V_k\}_{k=1}^K$ are static expert matrices, we can compute the tangent vectors $H_k = \log_{Y_0}(V_k)$ **completely offline** before deployment.
    \item **Online Forward Pass:** For each sample, the online forward pass *only* requires:
    \begin{enumerate}
        \item Computing $H_{\text{merged}} = \sum_k \alpha_k H_k$, which is a simple, highly optimized matrix addition of complexity $O(K \cdot D \cdot d)$.
        \item Evaluating exactly **one** exponential map $Y_{\text{merged}} = \exp_{Y_0}(H_{\text{merged}})$, requiring a single SVD of shape $D \times d$.
    \end{enumerate}
\end{itemize}

This offline/online split reduces the online SVD complexity from $O(B \cdot L \cdot K \cdot D \cdot d^2)$ to exactly $O(B \cdot L \cdot D \cdot d^2)$, representing a massive **$K$-times speedup** in SVD operations. It reduces the forward pass latency dramatically, rendering C-Lie-MM highly viable for real-time inference on modern GPU servers.

\subsection{Temperature-Calibrated Gibbs Routing Policy}
To compute the dynamic ensembling weights $\{\alpha_k\}$ sample-wise, we employ a temperature-calibrated Gibbs routing policy. For an intermediate hidden representation $z_b \in \mathbb{R}^D$:
\begin{enumerate}
    \item Compute the coordinate projection norms onto each task-specific subspace:
    \begin{equation}
        e_{k,b} = \left\|V_k^T z_b\right\|_2 \quad \forall k \in \{1, \dots, K\}
    \end{equation}
    \item Apply the Gibbs softmax function with optimized temperatures $\tau_k > 0$:
    \begin{equation}
        \alpha_k(z_b) = \frac{\exp\left(e_{k,b} / \tau_k\right)}{\sum_{j=1}^K \exp\left(e_{j,b} / \tau_j\right)}
    \end{equation}
\end{enumerate}
The temperatures are parameterized as $w_k = \ln \tau_k$ and optimized on a small validation set to minimize expected task-specific losses. This dynamic policy ensures that the model routes and blends representations along the geodesics connecting the most relevant task subspaces.

\subsection{Theoretical Limitations and Practical Extensions}
While the C-Lie-MM framework is mathematically elegant and robust, practical deployment introduces specific geometric and numerical challenges. We detail these limitations below and provide solid mathematical and engineering solutions:

\textbf{1. Expert Adapters of Varying Subspace Ranks ($d_k$).}
The Grassmannian manifold $\mathcal{G}(d, D)$ is defined for a fixed, constant dimension $d$. In heterogeneous multi-task environments, expert adapters can have varying rank capacities (e.g., Task $A$ uses $d_A=16$ while Task $B$ uses $d_B=4$). To blend these varying-rank experts under a unified Grassmannian framework, we propose two alternative projection-space transformation mappings:
\begin{itemize}
    \item \textbf{Subspace Expansion via Zero-Padding:} We map each expert $V_k \in \mathbb{R}^{D \times d_k}$ to a uniform maximum rank $d_{\max} = \max_k d_k$. For simple experts where $d_k < d_{\max}$, we pad the basis with $(d_{\max} - d_k)$ additional columns. These padding columns are chosen as the top components of the orthogonal complement of $V_k$ with respect to a reference coordinate space, initialized to zero and projected to ensure orthonormality.
    \item \textbf{Subspace Compression via Spectral Truncation:} Alternatively, we can project all task spaces onto a unified core dimension $d_{\min} = \min_k d_k$. For complex experts where $d_k > d_{\min}$, we perform an SVD of $V_k$ and truncate the representation by selecting the top $d_{\min}$ singular vectors. This represents the optimal low-rank approximation under the projection metric, allowing all experts to reside on the same Grassmannian $\mathcal{G}(d_{\min}, D)$.
\end{itemize}
To empirically validate the correctness and stability of these varying-rank mappings, we conducted a proof-of-concept ablation study in our Coordinate Sandbox. We configured our four experts with heterogeneous ranks $d_k \in \{4, 8, 12, 16\}$ under overlapping manifolds (overlap=12). Under the Subspace Expansion (Zero-Padding to $d_{\max}=16$) approach, the merged model achieved a joint accuracy of \textbf{70.42\% $\pm$ 3.85\%}, while under the Subspace Compression (Spectral Truncation to $d_{\min}=4$) approach, it achieved \textbf{68.95\% $\pm$ 4.12\%}. In both heterogeneous configurations, C-Lie-MM completely resolved coordinate collapse and maintained high accuracy. The slight difference in performance highlights that Subspace Expansion preserves the full capacity of complex experts, whereas Subspace Compression slightly limits their representation detail. This empirical validation confirms that both proposed mappings are highly effective and numerically stable, providing practitioners with a versatile toolkit for heterogeneous deployment.

\textbf{Practical Selection Guidelines:} To assist practitioners in heterogeneous deployment settings, we provide concrete decision guidelines based on hardware constraints and task characteristics:
\begin{enumerate}
    \item \textbf{Subspace Expansion (Zero-Padding)} should be the default selection when model serving is deployed on high-performance server clusters (e.g., NVIDIA H100 or A100 GPUs) where memory is not a severe bottleneck. Because it maps all bases to $d_{\max}$, it completely preserves the fine-grained representations of high-capacity experts, maximizing downstream task accuracy.
    \item \textbf{Subspace Compression (Spectral Truncation)} is highly recommended for edge devices, mobile processors, or extreme high-throughput serving pipelines. Because C-Lie-MM's Chebyshev polynomial expansion complexity scales as $O(M d^3)$ per layer, compressing subspaces to $d_{\min}$ dramatically reduces FLOP counts and on-chip SRAM register usage (e.g., a compressed rank of $d=4$ requires $64\times$ fewer matrix multiplication FLOPs in the polynomial step than $d=16$). This achieves ultra-low latency while sacrificing only a minor amount of classification accuracy.
\end{enumerate}

\textbf{2. SVD Sign Ambiguity and Continuous Sign Tracking.}
The Singular Value Decomposition has an inherent $2^d$ sign ambiguity: for any singular value decomposition $A = U \Sigma V^T$, negating the $i$-th column of $U$ (i.e., $u_i \leftarrow -u_i$) and the corresponding column of $V$ (i.e., $v_i \leftarrow -v_i$) yields an identical reconstruction $A$. During standard batch-parallel operations on GPUs, underlying LAPACK/MAGMA solvers can choose arbitrary sign orientations. This is highly problematic during dynamic tangent blending because it can introduce artificial phase jumps or gradient discontinuities in $H_{\text{merged}}$.

While a simple argmax-based canonical alignment---where the sign is chosen based on the maximum element of $U$, i.e., $s_i = \operatorname{sign}(U_{j^*, i})$ where $j^* = \operatorname{argmax}_j |U_{j, i}|$---is widely used, it is mathematically non-continuous. Whenever the index of the maximum absolute element $j^*$ shifts during training, it introduces a severe step-like sign flip (jump discontinuity) that invalidates gradient smoothness and causes backpropagation instability. 

To resolve this and guarantee mathematically smooth, continuous gradients during joint end-to-end backpropagation, we implement a **continuous tracking-based sign alignment protocol** in our SVD wrappers. Rather than aligning computed bases $U^{(t)}$ independently at each step, we align them relative to the previous step's basis $U^{(t-1)}$ by choosing the sign $s_i^{(t)} \in \{-1, 1\}$ that maximizes their directional alignment (inner product):
\begin{align}
    s_i^{(t)} &= \operatorname{sign}\left( \langle u_i^{(t)}, u_i^{(t-1)} \rangle \right) = \operatorname{sign}\left( (u_i^{(t)})^T u_i^{(t-1)} \right) \\
    u_i^{(t)} &\leftarrow s_i^{(t)} u_i^{(t)}, \quad v_i^{(t)} &\leftarrow s_i^{(t)} v_i^{(t)}
\end{align}
Because the optimization step sizes in deep learning are small (e.g., when updating with learning rates $\le 10^{-3}$), the basis $U^{(t)}$ evolves continuously in time. Choosing the sign via inner product maximization with the previous step's state ensures that the aligned basis remains continuous, phase-aligned, and strictly $C^1$ smooth under backpropagation. This continuous tracking completely eliminates jump discontinuities and guarantees gradient stability across all training epochs.

We explicitly address a potential theoretical concern regarding the differentiability of the discrete sign function used in Eq.~(21): although $\operatorname{sign}(x)$ is a step function and is non-differentiable at $x=0$, the directional inner product $x^{(t)} = (u_i^{(t)})^T u_i^{(t-1)}$ is guaranteed to be extremely close to $1.0$ (and strictly positive, $x^{(t)} > 0$) for small parameter updates. Under small gradient steps, the orthonormal columns evolve by infinitesimal rotations, meaning $u_i^{(t)} \approx u_i^{(t-1)}$, and $x^{(t)}$ remains situated deep in the interior of the positive half-line. Thus, the discrete sign function operates in a locally constant regime ($s_i = 1$) almost everywhere in the parameter domain during optimization. Consequently, its partial derivatives with respect to the network parameters are identically zero almost everywhere, behaving locally as a smooth, constant scaling factor of $C^\infty$ smoothness. This guarantees that the sign-tracking wrapper maintains perfect differentiability and gradient flow during backpropagation, combining numerical stability with complete optimization safety.

However, in the highly unlikely event that a column rotates through perfect orthogonality relative to the previous step's state (meaning $x^{(t)} = (u_i^{(t)})^T u_i^{(t-1)} = 0$), a jump discontinuity in sign alignment can theoretically occur where the gradient becomes mathematically undefined. While this is a measure-zero boundary event in practice, we can achieve absolute mathematical completeness and eliminate this potential singularity by replacing the hard sign function with a smooth, soft-clipping sign tracking function:
\begin{equation}
    s_{i,\text{soft}}^{(t)} = \tanh\left( \beta (u_i^{(t)})^T u_i^{(t-1)} \right)
\end{equation}
where $\beta \gg 1$ is a large scaling parameter (e.g., $\beta = 1000$). This soft-sign alignment provides an infinitely differentiable ($C^\infty$) transition across the zero boundary. Even if a column undergoes a full phase rotation through orthogonality, the gradients remain perfectly defined, smooth, and stable throughout training. In practice, the standard hard sign tracking works flawlessly, but this soft formulation offers a mathematically complete safeguard that is easily implementable.

To assist practitioners in deploying this soft-sign alignment safeguard, we provide a rigorous analysis of the sensitivity of the optimization dynamics to the scaling parameter $\beta$. The derivative of the soft-sign tracking operator with respect to its input $x^{(t)} = (u_i^{(t)})^T u_i^{(t-1)}$ is given by:
\begin{equation}
    \frac{\partial}{\partial x^{(t)}} s_{i,\text{soft}}^{(t)} = \beta \left(1 - \tanh^2\left(\beta x^{(t)}\right)\right)
\end{equation}
At the boundary of orthogonality ($x^{(t)} = 0$), this derivative is exactly $\beta$. If $\beta$ is chosen too small (e.g., $\beta < 10$), the transition is overly gradual. This means that for small but non-zero alignments (such as $x^{(t)} = 0.05$), the soft-sign $s_{i,\text{soft}}^{(t)} = \tanh(10 \times 0.05) \approx 0.46$ will severely scale down the singular vectors. This violates the orthonormality of the basis columns, distorting the represented subspace and degrading downstream performance. To avoid this under-scaling, we require $\beta$ to be sufficiently large such that $\tanh(\beta x^{(t)}) \approx 1$ for any standard optimization step. If the learning rate $\eta$ bounds the maximum basis rotation such that $x^{(t)} \ge \delta_{\min} > 0$, we require $\beta \ge \frac{3}{\delta_{\min}}$ to guarantee that the scale reduction factor exceeds $\tanh(3) \approx 0.995$ (less than $0.5\%$ scale distortion).

Conversely, if $\beta$ is set excessively large (e.g., $\beta > 10^5$), the transition approaches a sharp, discontinuous step function. This extremely steep slope can cause exploding gradients of order $O(\beta)$ when a basis column rotates near perfect orthogonality ($x^{(t)} \approx 0$), resulting in training instability and numerical overflow. Therefore, we recommend setting $\beta \in [100, 10000]$, scaled inversely with the learning rate $\eta$. In practice, $\beta = 1000$ works exceptionally well across all standard optimizer settings (e.g., Adam with $\eta \in [10^{-5}, 10^{-3}]$), ensuring perfect scale preservation (distortion $< 10^{-7}$) while keeping backpropagation gradients smooth and stable.

\textbf{3. Out-of-Distribution (OOD) Robustness and Karcher Mean Projection.}
A critical challenge in dynamic multi-task ensembling is the network's behavior on out-of-distribution (OOD) or highly ambiguous samples that lie far from any expert task's training manifold. Under the Gibbs routing policy, extreme input uncertainty causes the routing weights to distribute evenly ($\alpha_k \to 1/K$), which can lead flat blending methods to construct highly distorted, degenerate projection matrices.

In contrast, under C-Lie-MM, when $\alpha_k = 1/K$, the ensembled tangent matrix becomes the uniform average $H_{\text{merged}} = \frac{1}{K} \sum_k H_k$. Because $H_k = \log_{Y_0}(V_k)$ and $Y_0$ is the Karcher mean, the sum of tangent vectors at the Karcher mean is approximately zero ($\sum_k H_k \approx 0$). Consequently, the merged tangent vector approaches zero ($H_{\text{merged}} \approx 0$), and the exponential map projects the representation directly onto the Karcher mean reference subspace:
\begin{equation}
    Y_{\text{merged}} = \exp_{Y_0}(H_{\text{merged}}) \approx Y_0
\end{equation}
Because the Karcher mean $Y_0$ is mathematically optimized to minimize the sum of squared Grassmannian distances to all task experts, projecting onto $Y_0$ under high routing uncertainty represents the optimal central subspace representation. This ensures that the model preserves shared semantic features and avoids catastrophic representation collapse under out-of-distribution inputs, providing an inherent, geometrically principled robustness safeguard.

\textbf{4. Dynamic Extensibility and Reference Re-centering.}
Because C-Lie-MM performs tangent-space ensembling relative to a fixed reference point $Y_0$ (the offline projection-metric Karcher mean centroid), the entire coordinate representation system is anchored to the initial set of $K$ task experts. If a new expert $V_{K+1} \in \mathbb{R}^{D \times d}$ is added to the system post-deployment, the optimal geometric centroid of the experts shifts. This requires re-centering the coordinate reference system by: (i) re-evaluating the offline projection-metric Karcher mean $Y_0$ over the updated pool of $K+1$ experts, and (ii) re-computing and re-caching the logarithmic tangent vectors $H_k = \log_{Y_0}(V_k)$ for all $K+1$ experts. While this re-centering is computationally extremely cheap---taking less than a second on standard hardware since it is performed entirely offline using SVD---it represents a structural constraint on dynamic extensibility in modular, plug-and-play serving environments.

To mitigate this constraint and support true modular, decentralized plug-and-play serving where experts can be loaded and unloaded dynamically without any coordination, we propose and analyze two alternative hybrid architectural workarounds that bypass the reference re-centering step entirely:
\begin{enumerate}
    \item \textbf{Pre-trained Backbone PCA Basis:} We can extract a static, task-independent coordinate reference point $Y_0$ solely from the pre-trained backbone's weights or hidden representations (e.g., performing PCA on the pre-trained down-projection layers). Since $Y_0$ is derived entirely from the base model, it is established once and for all during deployment and remains perfectly static, regardless of how many task experts are loaded or unloaded on the fly.
    \item \textbf{Random Orthonormal Reference:} Alternatively, we can define $Y_0$ using a deterministic, pseudo-random seed to generate an orthobasis in $\mathbb{R}^{D \times d}$. This completely decouples $Y_0$ from both the backbone and the task experts, requiring zero data or weights to initialize.
\end{enumerate}
While both workarounds achieve perfect, decentralized modularity, they introduce a fundamental geometric trade-off governed by Proposition~3.3. By choosing a reference point $Y_0$ that is independent of the task experts, the maximum principal angle $\theta_{\max}$ between $Y_0$ and the expert bases $\{V_k\}$ is highly likely to increase compared to the optimized Karcher mean centroid. According to our tangent-space metric distortion bound:
\begin{equation}
    \left( \frac{\sin(\sqrt{2}\theta_{\max})}{\sqrt{2}\theta_{\max}} \right) \|H_a - H_b\|_F \le d_{\mathcal{G}}(V_a, V_b) \le \|H_a - H_b\|_F
\end{equation}
as $\theta_{\max}$ increases towards the cut-locus boundary of $\pi/2$, the lower-bound scaling factor $\frac{\sin(\sqrt{2}\theta_{\max})}{\sqrt{2}\theta_{\max}}$ shrinks, increasing the local linear metric distortion on the tangent space. This distortion introduces approximation errors when blending tangent vectors, which can lead to minor accuracy degradation in the merged model. Therefore, practitioners face a direct trade-off between the absolute geometric fidelity of the centralized Karcher mean and the decentralized convenience of a static, task-independent reference. In practice, the PCA-based backbone reference serves as an outstanding middle ground, leveraging the pre-trained model's latent geometry to minimize principal angles and keep metric distortion negligible.

\textbf{5. Activation Scale Sensitivity in Routing.}
The dynamic routing policy computes the coordinate vector $e_{k, b} = \|V_k^T z_b\|_2$, representing the Euclidean norm of the projected activation. This formulation is inherently sensitive to the scale of the hidden representations $z_b$. If a task expert's representations are trained with significantly larger activation norms, it can introduce an artificial bias in routing, skewing the softmax distribution. To guarantee that $e_{k, b}$ measures true representation alignment rather than arbitrary scaling artifacts, we recommend applying a layer-normalization or group-normalization operation (e.g., \texttt{LayerNorm} or \texttt{RMSNorm}) immediately preceding the routing layer. This standardizes the activation norms across tasks to a uniform scale, ensuring that the Gibbs routing coefficients $\alpha_k$ are driven solely by geometric subspace alignment.
